\documentclass[../main.tex]{subfiles}
\begin{document}

\chapter{Transitorio Capacitivo Resistivo}
\img{images/09/transitorio}{Transitorio capacitivo resistivo.}
\info{Transitorio Capacitivo-Resistivo}{ \emph{Cuando se produce un cambio en las magnitudes de un circuito, tensión o corriente, decimos que el circuito está en régimen transitorio. Al cambiar las condiciones de un elemento de un circuito se pierde el régimen permanente, y tras
sucederse los cambios de tensión y corriente se vuelve de nuevo al
equilibrio en otro régimen permanente. Al intervalo entre los dos regimenes permanentes se le denomina régimen transitorio.}}

\info{Constante de tiempo}{\begin{equation*}
    \tau = RC
\end{equation*} $\tau$ se llama constante de tiempo del circuito. Es un indicador de la velocidad de reacción del circuito ante una perturbación (debido a un escalón de tensión). Cuanto mayor sea este valor, el valor final del estado de equilibrio se alcanzará más rápidamente.}
\begin{figure}[h!]
    \centering
    \begin{tikzpicture}
    \draw (0,0)--(0,4)node [left]{$I$} [-stealth];
    \draw (0,0)--(6,0)node [below]{$t$} [-stealth];
    \draw[scale=1, domain=0:5, smooth, variable=\x, gray] plot ({\x}, {2.71828*2.71828^(-\x)});
    \draw [dashed](1,0)node [below]{$\tau$}--(1,1)[-stealth];
    \draw [dashed](0,1)node [left]{$37\%I_0$}--(1,1);
    \draw (0,2.7) node [left]{$I_0$};
    \end{tikzpicture}
    \caption{Curva de descarga de un capacitor}
    \label{fig:my_label}
\end{figure}

\actividad{Resolver los siguientes ejercicios de transitorio eléctrico}

\begin{enumerate}
    \item Encontrar la constante de tiempo $\tau$ para un circuito formado por un capacitor de $C=10\mu F$ conectado con una resistencia de $R=10k\Omega$.
    \item Si la constante de tiempo es $\tau=1s$ y se dispone de una resistencia de $R=100k\Omega$, encontrar la capacidad del capacitor $C$ en Faradios.
    \item Encontrar cuanto tiempo llevará que se cargue un capacitor de $200\mu F$ al 63\% si se conecta a una resistencia de $R=100\Omega$. 
    \item Encontrar cuanto tiempo llevará que se cargue un capacitor de $200\mu F$ al 63\% si conecta a 220V y la corriente inicial es de 2A. 
    \item cuanto tiempo demorará en descargarse al 37\% Un capacitor de $C=1F$ que fue cargado con una corriente inicial $I_0=10A$ y $V=24V$
    
    \info{Carga completa de capacitor 99\%}{
    Un capacitor se carga al 99\% después de 
    $t=5\tau$}
    
    \item  Encontrar el tiempo $t$ que demorará en cargarse al 99\% un capacitor de $C=10\mu F$ conectado a una resistencia $R=1k\Omega$. 
    
    \item Encontrar la constante de tiempo $\tau$ de los siguientes circuitos entre $a$ y $b$: 
    \begin{enumerate}
        \item . 
            \begin{center}
        \begin{circuitikz}
         \draw (-3,0) node[label={above:a}]{} to[resistor=$3k\Omega$,*-*](2,0)  (2,0)--(2,2) to[resistor=$1\Omega$] (-2,2)--(-2,0);
         \draw (-2,0)--(-2,-2);
         \draw (2,0)--(2,-2);
         \draw (-2,-2) to[resistor=$2k\Omega$] (2,-2);
         \draw (2,0) to[capacitor=$100\mu F$] (4,0);
         \draw (4,0) to[resistor=$100\Omega$,-*] (6,0) node[label={above:b}]{};
        \end{circuitikz}
        \end{center}
    \item .
            \begin{center}
        \begin{circuitikz}
         \draw (-3,0) node[label={above:a}]{} to[capacitor=$3\mu F$,*-*](2,0)  (2,0)--(2,2) to[capacitor=$1\mu F$] (-2,2)--(-2,0);
         \draw (-2,0)--(-2,-2);
         \draw (2,0)--(2,-2);
         \draw (-2,-2) to[capacitor=$2\mu F$] (2,-2);
         \draw (2,0) to[capacitor=$5\mu F$] (4,0);
         \draw (4,0) to[resistor=$10\Omega$,-*] (6,0) node[label={above:b}]{};
        \end{circuitikz}
        \end{center}
    
       \item .
        \begin{center}
        \begin{circuitikz}
         \draw (-3,0) node[label={above:a}]{} to[resistor=$10\Omega$,*-](3,0) to[capacitor=$200\mu F$,-*] (3,-3) node[label={right:b}]{};
        \draw (-3,0)  to[resistor=$200\Omega$,*-](-3,-3) to[resistor=$1\Omega$,-*] (3,0);
        \end{circuitikz}
        \end{center}
        
               \item .
        \begin{center}
        \begin{circuitikz}
         \draw (-3,0) node[label={above:a}]{} to[capacitor=$10\mu F$,*-](3,0) to[resistor=$200\Omega$,-*] (3,-3) node[label={right:b}]{};
        \draw (-3,0)  to[capacitor=$200 \mu F$,*-](-3,-3) to[capacitor=$1\mu F$,-*] (3,0);
        \end{circuitikz}
        \end{center}
        
        \item .
              \begin{center}
        \begin{circuitikz}
           \draw (0,0) node[label={left:a}]{} to[capacitor=$1\mu F$,*-*] (4,0) 
           to[resistor=$6\Omega$,*-*] (7,2)--(7,5)node[label={right:b}]{} 
           to[resistor=$2\Omega$,*-*] (3,5)--(3,2)
           to[resistor=$1\Omega$,*-*] (1.5,1)--(0,0);
           \draw (0,0) to[capacitor=$6\mu F$,*-*] (0,4)
           to[capacitor=$3\mu F$,*-*] (4,4)--(4,0);
           ; 
        \end{circuitikz}
        \end{center}    
    \end{enumerate}
    
    
    
    
    
\end{enumerate}






\end{document}